% !TEX root = ../main.tex
\section{Objections and Replies}
\label{sec:objections}

\subsection{Interpretationism}
\label{sec:interpretationism}

\paragraph{Objection.} On an interpretationist view of the attitudes
\citep{dennett1987}, to have beliefs just is to be a system whose behaviour is
voluminously and reliably predictable from the intentional stance---by ascribing
beliefs and desires and assuming rationality. Language models are prime cases.
Users predict their outputs by asking what something that believed its training
data would say, and the predictions succeed at scale. By the interpretationist's
own standard, then, these systems are believers, and the formal profile of
$\Out_M$ is beside the point: human believers do not validate KD45 either, and
Dennett never held that the reality of belief required it.

\paragraph{Reply.} The objection understates what the intentional stance
requires. Dennett's criterion is not predictability by any mentalistic heuristic.
It is predictability \emph{under an assumption of rationality}: the ascribed
beliefs must be treated as roughly closed under obvious consequence, roughly
consistent, and coupled to desires through practical inference. That assumption
is what Section~\ref{sec:output-operator} shows to fail, not at the margins but
as the central phenomenon.

The user who predicts a model well is not working from the intentional stance.
She is working from what we might call the \emph{corpus stance}---what would text
shaped like the training distribution look like here?---supplemented by knowledge
of prompt-form sensitivity. Anyone who has spent time getting useful output from
these systems will recognise the second half as the larger part of the skill, and
will recognise that it is not a skill one needs for predicting believers.

The two stances make divergent predictions, and they diverge precisely at the
RE-failures. The intentional stance predicts stability under paraphrase: a
believer asserts the same proposition however you ask. The corpus stance predicts
formulation-sensitivity. The corpus stance wins those contests systematically. By
Dennett's own lights, a system whose behaviour is better predicted from a
design-level stance than from the intentional stance is one for which intentional
ascription is dispensable metaphor. Interpretationism, taken seriously, is an ally
of my conclusion rather than an objection to it.

Note that this reply does not require the strong claim rejected in
Section~\ref{sec:internal}. It is consistent with there being
prompt-invariant representations inside. What it denies is that belief
attributions are \emph{projectible}---that they support prediction across the
transformations under which a believer's attitudes are stable.

There is a residual instrumentalist fallback: even if the corpus stance predicts
better, belief-talk is a convenient compression. Granted, as loose talk. But
Section~\ref{sec:hallucination} showed the compression is not harmless. It
imports truth-conditional evaluation and person-level appraisal, both of which
misclassify the phenomena that matter. Where a metaphor systematically distorts
measurement, instrumentalism recommends retiring it.

\subsection{The human-parallel objection}
\label{sec:human-parallel}

\paragraph{Objection.} This is the deepest challenge, and it deserves the most
space. Humans, the objection runs, exhibit every failure catalogued in
Section~\ref{sec:output-operator}. Framing effects
\citep{tversky1981} are RE-failures: subjects accept a treatment described as
ninety per cent survival and reject the same treatment described as ten per cent
mortality. Closure failures are K-failures. The conjunction fallacy is worse than
an M-failure, since subjects rank a conjunction above its conjunct. Inconsistency
is a D-failure. Confabulation research shows 4-failures: people's reports of their
own attitudes are generated by processes insensitive to those attitudes. And
negative introspection fails whenever we are wrong about what we do not believe,
which is often.

If these failures showed that $\Out_M$ is not a belief operator, parity of
reasoning shows human belief is not one either. So either the argument proves too
much---nothing believes anything, and the contrast the paper needs
evaporates---or the normal-operator standard is the wrong standard, in which case
premise (1) of Section~\ref{sec:structural-argument} fails. Worse: the human data
suggest belief was never well modelled by KD45, so the received formal sense is a
standard nothing meets, and disqualifying machines by it is empty.

\paragraph{Reply.} The objection trades on treating all deviations from normality
alike. The competence/performance distinction breaks the parity, but it has to be
deployed carefully, because appeals to competence can degenerate into the
unfalsifiable claim that humans really respect logic deep down. The claim I need
is more disciplined: for humans, closure-conformity is the \emph{functional norm}
of the belief system, in the sense that the system contains mechanisms whose
function is to detect and repair closure violations, and whose operation explains
characteristic patterns in the data. Three such patterns, none of which $\Out_M$
exhibits.

\paragraph{First: normative acknowledgement.} When a framing effect or
conjunction fallacy is exposed to a human subject---when the equivalence of the
formulations is made salient---the subject recognises the earlier response as an
error \emph{by her own standard}. She does not need to be persuaded that
consistency was required; she already held that it was. Whether debiasing
reliably produces lasting behavioural transfer is a contested empirical question
and I do not rest anything on it. What matters is the acknowledgement, because
acknowledgement is what reveals a standing internal standard relative to which
the violation registers as a violation.

There is no analogue here. Correcting a model within a context window produces
local, context-bound compliance that evaporates with the context; the underlying
assertion dispositions across prompt-forms are untouched. Nothing internal to the
system treats the RE-violation as a violation. In humans the framing effect is a
bug relative to the system's own error-correcting standard. In these systems
formulation-sensitivity is not a bug relative to any standard, because there is no
standard---it is the mechanism.

\paragraph{Second: resource-sensitivity.} Human closure failures are modulated by
the variables that indicate a bounded competence. They diminish with attention,
incentive, time, expertise, and cognitive unloading, and increase under load.
That is the fingerprint of a closure-respecting process running out of resources.
The corresponding question for machines is whether their failures show the same
modulation, and it is a real question rather than a rhetorical one. I take it up
separately in Section~\ref{sec:reasoning}, because the systems that bear on it
most directly are recent.

\paragraph{Third: the cross-context practical-reasoning economy.} This is the
consideration I want to press hardest, because it explains the other two rather
than sitting beside them.

Human belief is embedded in an economy that machine assertion-states do not
participate in. A human's belief that $\phi$ is poised, in principle, to combine
with any of her desires and any of her other beliefs to shape action across
arbitrary contexts. The belief that a bridge is unsafe interacts with the desire
to survive, the plan to travel, the testimony offered to others, the wagers
accepted, the warnings given to strangers. It is this cross-context inferential
and practical promiscuity that makes closure a constitutive norm for belief rather
than an optional refinement. A state that will be recruited by open-endedly many
projects must be kept consistent and closed, on pain of the projects failing, and
the pressure of that requirement is what the repair mechanisms of the first
pattern subserve. Belief is closure-normed because it is practically integrated,
and it is practically integrated because the agent bears the costs.

Machine assertion-states have no such economy. There is no standing goal-structure
they serve, no cross-context recruitment---each assertion is minted in its context
and expires with it---and hence no functional pressure toward closure and nothing
whose job is to enforce it. The absence of the practical-reasoning economy is not
one more disanalogy on a list. It is why the disanalogy pattern takes the shape it
does. Humans fail closure while being governed by it; these systems fail closure
while being governed by nothing.

So the parity dissolves. Human belief is a normal operator at the level of
competence, with non-normal performance; the KD45 idealisation is an idealisation
of something---the functional norm the repair machinery enforces. $\Out_M$ is
non-normal at the level of competence, and there is nothing for KD45 to idealise.
The human data do not show that belief was never KD45-like. They show that KD45
describes belief's regulative standard rather than its statistical surface, which
is what an idealisation was always supposed to do. And the objection's disjunction
is a false dilemma: the standard can be a competence-level standard that humans
meet and these systems do not.

A residual worry deserves acknowledgement. The competence/performance distinction
is itself contested, and one might fear it functions here as an immunising
stratagem---any human failure filed under performance, any machine failure under
competence, the game rigged from the start. The reply is that the filing is done
by independent and empirically checkable criteria, and each could come out the
other way. A system exhibiting normative acknowledgement that persisted across
prompt-forms, a resource-sensitive rather than frequency-sensitive failure
profile, and a persistent cross-context economy recruiting its assertion-states
into open-ended practical reasoning would satisfy the marks I have used to
separate the cases, and I would owe the objection its conclusion. That is not a
concession made for form's sake; the second criterion is under active pressure
right now.

\subsection{Reasoning models and the resource-sensitivity criterion}
\label{sec:reasoning}

The second signature is the one most exposed to recent developments, and I would
rather meet it directly than let the scope note do the work.

Models trained with reinforcement learning on verifiable outcomes improve on
closure-heavy tasks---mathematics, formal deduction, multi-step
inference---as a function of how much computation they spend at inference time.
Extended deliberation is a plausible machine analogue of the pencil and paper I
cited above as a human competence-signature. If additional resources reliably buy
better closure, the modulation profile looks resource-sensitive, and the second
signature stops separating the cases.

The reply is to be precise about what the criterion measures. My thesis concerns
congruentiality, not task accuracy, and these are different quantities. The
question is not whether extra computation raises benchmark scores. It is whether
extra computation makes assertion \emph{invariant under logically equivalent
paraphrase}. A system can improve on both counts, on one, or on neither, and the
framework makes a definite prediction about which.

Two outcomes are possible and they mean different things.

If paraphrase-divergence falls as computation rises---if the system converges on
the same answer across contraposed, negation-dual, and re-notated formulations
that it previously split on---then the system is moving toward classicality. The
framework says so without embarrassment. Section~\ref{sec:intro} was explicit
that nothing here rules out architectures whose assertion operators approach
normality, and Section~\ref{sec:checklist} states what that would require. A
result of this kind would be the framework working, not failing.

If accuracy rises while paraphrase-divergence holds steady, then the extra
computation is broadening the set of formulations the system handles without
improving the operator's relation to propositional content. That is coverage, not
competence, and it is what the architectural argument of
Section~\ref{sec:re-failure} predicts: more compute routes generation through more
of the formulation space that training covers well, and leaves the individuation
of contents by strings exactly where it was.

The discriminating experiment is therefore not hard to describe, and
Section~\ref{sec:re-audit} describes it: run an equivalence battery across
inference-compute budgets and across training checkpoints, and watch whether the
divergence statistic falls or merely shifts. I am not aware that this measurement
has been made systematically, which is itself worth saying. The claim that
reasoning training produces closure competence rather than broader coverage is,
at present, an assumption on both sides of the argument.

One further consideration supports my side of it without settling the matter.
Where the reasoning trace has been examined, it is frequently not the causal route
to the answer: models produce rationales that do not reflect the computation that
produced the output. If the trace is post hoc, then whatever the additional
computation is buying, it is not the surveyable inferential closure the human
competence-signature consists in.
% TODO(cite): chain-of-thought unfaithfulness. I could not verify the Turpin et
% al. reference; verify before relying on it, since it carries weight here.

\subsection{Deflationism about the terminological question}
\label{sec:deflationism}

\paragraph{Objection.} Suppose everything above is right. Why does it matter
whether we call the underlying states beliefs? The word is ours to extend.
Natural-kind terms routinely stretch to cover new cases---``memory'' now covers
RAM, ``computation'' covers brains---and no one thinks these extensions embody
metaphysical errors. Perhaps ``belief'' is polysemous, or a cluster concept, or
should be precisified by stipulation as engineering convenience dictates. The
grand negative thesis then shrinks to a lexicographical preference.

\paragraph{Reply.} The terminological question is substantive, and it is
substantive for a reason that has nothing to do with semantic conservatism.
``Belief'' is not a bare classifier. It is a node in an inferential and practical
web, and calling a state a belief licenses a bundle of further moves: expecting
stability under paraphrase, holding the system to consistency, reading its
self-reports as introspective, applying truth-conditional success criteria,
extending trust on the basis of past accuracy, and assigning the kind of appraisal
that presupposes a subject. Every element of that bundle misfires here, and
Sections~\ref{sec:hallucination} and~\ref{sec:evaluation} show the misfires have
measurable costs: benchmarks that reward cell-(ii) luck, users who extrapolate
reliability across prompt-forms where none transfers, procurement and regulatory
language that asks whether a system ``knew'' a fact.

So the case against extending ``belief'' is that this particular extension exports
a false theory. Where ``memory'' for RAM exported a functional abstraction whose
distortions have been comparatively mild---though the storage-and-retrieval
picture of human memory has itself drawn criticism, so the contrast case is not
pristine---``belief'' for $\Out_M$ exports normality assumptions the system
provably violates. Terminological choices are cheap only when the inferential
freight is cheap. Here it is not.

The cluster-concept variant is the strongest form of the objection, and the
apparatus of Section~\ref{sec:monadic} answers it. Grant that ``belief'' is a
cluster concept. The question is then which members of the cluster are central,
and the answer is given by what the formal tradition idealises: the
closure-related properties, which is to say congruentiality, consistency, and
closure under obvious consequence. Conditions (1) and (2) of the monadic
core---informational, behaviour-guiding---are peripheral members, and I grant them
freely. Condition (3) denies the central ones. A cluster concept whose central
members are absent is not satisfied, and the sub-doxastic category exists exactly
to name what is left.

Finally, a concession that costs nothing. The formal and methodological results of
this paper are independent of the label. Machine assertion is a hyperintensional
non-normal operator; it therefore admits no Kripke semantics and no bare
neighborhood semantics; grounding cross-cuts truth; accuracy bounds nothing about
the grounding rate; and hallucination so defined requires attribution audit to
measure. Call the states ``beliefs*'' if you like and every one of those claims
stands. But it does not follow that the doxastic question was idle, because the
argument just given shows the label carries commitments the results contradict.
Deflationism about the word is not deflationism about the theory the word smuggles
in.
